% !TeX spellcheck = ru_RU
% !TEX root = vkr.tex

\section{Угрозы валидности} \label{sec:threats}

В настоящем разделе обсуждаются факторы, способные повлиять на интерпретацию полученных результатов. Обсуждение структурировано по четырём классическим категориям: внутренняя, внешняя, конструктная и статистическая валидность.

\subsection{Внутренняя валидность}

\textbf{Один прогон инференса.} Каждая модель оценена на каждом бенчмарке в одном прогоне без усреднения. Сэмплирование при \texttt{temperature}~$= 0{,}1$ не является строго детерминированным, что вносит остаточную стохастичность в число ошибок на уровне отдельных единиц. Эффект консервации ошибок (раздел~\ref{subsec:error_analysis}), зафиксированный при воспроизведении запусков, указывает на то, что большая часть ошибок систематична; однако для моделей, различия между которыми близки к ошибке измерения (например, \texttt{lite\_latest\_v2} и \texttt{lite\_rc\_v1} с $p = 0{,}73$), вывод об эквивалентности следует принимать с учётом этой неопределённости.

\textbf{Неуправляемые гиперпараметры LoRA.} Ранг $r$, коэффициент $\alpha$, learning rate и расписание обучения контролируются платформой Yandex Cloud и не доступны напрямую. Следовательно, делать выводы о влиянии конкретных гиперпараметров LoRA на результат нельзя; наблюдаемые эффекты отражают совокупный результат <<данные~$\times$~платформенная конфигурация>>.

\textbf{Субъективность разметки.} Первичная разметка выполнена одним автором; спорные случаи обсуждались с двумя помощниками. Формальное межэкспертное согласие (Cohen's~$\kappa$) не вычислялось. Это означает, что часть ошибок моделей на пограничных запросах (см.~<<трудное ядро>> в разделе~\ref{subsec:error_analysis}) может отражать неоднозначность самой разметки.

\subsection{Внешняя валидность}

\textbf{Одно семейство моделей.} Все эксперименты проведены на YandexGPT~5.1 (Lite и Pro). Переносимость выводов на модели других архитектур и провайдеров (GigaChat, GPT-4o, Llama, Qwen) не гарантируется.

\textbf{Узкий домен и язык.} Результаты получены на русскоязычных запросах в домене поиска недвижимости. Обобщение на другие домены (медицина, юриспруденция, e-commerce) требует отдельной валидации, особенно в части <<трудного ядра>>, которое в нашем случае состоит преимущественно из топонимов-аббревиатур и сложных составных запросов.

\textbf{Бинарность постановки.} В работе рассматривается выбор между двумя инструментами. При $|\mathcal{T}| > 2$ задача усложняется (растёт число возможных ошибок, снижается эффективная точность случайного угадывания), и количественные выводы~--- прежде всего минимальный объём обучающих примеров~--- не обязаны сохраняться.

\subsection{Конструктная валидность}

\textbf{Что именно измеряется.} Error rate и Macro~$F_1$ вычислены по результату классификации \emph{первого} вызова инструмента (\texttt{tool\_call}) в ответе модели. Качество последующего многошагового взаимодействия (передача корректных аргументов, обработка результата инструмента, диалоговая связность) не оценивается. Низкое значение error rate гарантирует корректный выбор функции, но не гарантирует корректность её применения.

\textbf{Метрика <<правильности>>.} Ошибка I и II рода определены в терминах классификации (раздел~\ref{sec:background}), а не в терминах удовлетворённости пользователя. В промышленной системе цена ошибки типа TAG$\to$GEO и GEO$\to$TAG может быть различна; данное различие в настоящей работе не учитывается.

\subsection{Статистическая валидность}

\textbf{Ограниченный объём выборок.} Бенчмарки содержат 500 и 250~примеров; для оценки модели с error rate $\leqslant 0{,}01$ это означает всего несколько ошибок в абсолютном выражении. Попарные сравнения на таких уровнях точности близки к границе статистической мощности, что ограничивает способность обнаруживать малые, но реальные различия между топ-моделями.

\textbf{Консервативность поправки Бонферрони.} При $\binom{12}{2} = 66$ попарных сравнениях поправка Бонферрони обеспечивает FWER~$\leqslant 0{,}05$, но предположение о независимости тестов нарушается (оценки моделей коррелируют, поскольку вычислены на общей выборке). Как следствие, фактический FWER ниже номинального, а мощность каждого отдельного теста снижается. Альтернативные методы (например, процедура Хохберга) могли бы дать более мощное, но менее консервативное решение; в настоящей работе мы намеренно выбрали консервативный вариант.

\textbf{Совпадение тренировочного и тестового распределений для \texttt{train\_real\_250}.} И \texttt{benchmark-real}, и \texttt{train\_real\_250} сформированы из одного и того же пула 500~реальных запросов (раздел~\ref{subsubsec:data_collection}). При отсутствии пересечений на уровне индивидуальных запросов (что проверено) распределения признаков всё равно близки. Поэтому высокие метрики \texttt{train\_real\_250} на \texttt{benchmark-real} отражают не только качество дообучения, но и распределенческое совпадение; выводы о превосходстве реальных данных над синтетическими следует формулировать именно в таком виде.
